\lesson{6}{Oct 1 2021 Fri (9:29:01)}{The Bill of Rights}{Unit 2}

The Bill of Rights is the name given to the first 10 amendments to the Constitution. It contains a list of individual rights that Americans have that the government must respect. Early leaders of the U.S. wanted to limit government so that the rights found in the \textbf{Bill of Rights} could not be taken away. The preamble to the \textbf{Bill of Rights} explains its purpose:

\begin{figure}[h]
    \begin{quote}
        \it{"The Conventions of a number of the States, having at the time of their adopting the Constitution, expressed a desire, in order to prevent misconstruction or abuse of its powers, that further declaratory and restrictive clauses should be added: And as extending the ground of public confidence in the Government, will best ensure the beneficent ends of its institution."}
    \end{quote}
    \caption{Preamble to the Bill of Rights}
\end{figure}

The \textbf{First Amendment} is probably the most famous. This amendment protects a person's rights to:

\begin{itemize}
    \item Free speech
    \item Press
    \item Assembly
    \item Petition
    \item Religious expression
\end{itemize}

Yet these terms can be difficult to define and apply to specific situations. Federal courts, including the U.S. Supreme Court, have interpreted the meaning of the amendments and the Constitution, carefully explaining how their judgments apply to the case. The justices base their rulings on one or more of these amendments, other laws, and precedent cases.

\subsubsection*{Absolute Individual Rights}

The broad language of the \bf{Constitution} and the \bf{Bill of Rights} could lead judges to interpret almost any action as protected. If a group of teens message each other about plans for a violent attack on another teen, should authorities ignore the messages because the teens have the right to \it{"Free Speech"}? You would not want such social network messages ignored, especially if you were the target of the plans.

Individual rights are not absolute. If everyone could do what they wanted, the rights of others would be violated and ultimately lead to chaos. Government must balance individual rights with the public good or general welfare of the people. While the language of the \bf{Constitution} and amendments protect individual rights, it also gives government the power to pass laws to restrict behavior, as stated in the \bf{Constitution's Preamble}, to \it{"establish Justice, insure domestic tranquility, provide for the common defense, (and) promote the general welfare"} of the country.

Legislation, the \bf{Bill of Rights}, court decisions, and the amendment process work together to both safeguard individual rights and limit them as deemed appropriate to protect the rights to safety and security for all. When people bring a case to the courts, judges can decide whether or not the Constitution allows a law that restricts the rights of individuals.

\paragraph{Amendment I}

The \bf{First Amendment} protects different aspects of \it{“Free Expression”}. People may share opinions, meet in groups, and publish their ideas in newspapers and on the Internet. They have the right to choose what religion to abide or to choose not to follow one. Americans do not have to worry about arrest for speaking about their beliefs. 
Yet they do often disagree on what constitutes \it{“Free Speech”}. A common debate centers on television violence and obscenity. Should the government censor programming that some Americans believe amoral but others consider to have artistic merit?

Laws and the courts have limited speech that threatens public safety or someone else's rights. For example, imagine someone yells "Shark!" as a practical joke at a crowded beach. In the rush to leave the beach, some people are hurt. Another example is if one teen spreads false stories about another on the Internet. Slander means lies spoken about a person that harms their reputation, and libel is the printed version of the same. These people can get in trouble with the law, although it is difficult to prove intentional slander or libel in court. An example related to assembly is that cities can require permits to have a festival in the park in order to plan extra police monitoring to ensure safety. In general, laws aimed at protecting public safety and preventing harassment limit the right to free speech.

\begin{figure}[h]
    \begin{quote}
        \it{“Congress shall make no law respecting an establishment of religion, or prohibiting the free exercise thereof; or abridging the freedom of speech, or of the press; or the right of the people peaceably to assemble, and to petition the Government for a redress of grievances.”}
    \end{quote}
    \caption{Amendment I}
\end{figure}

\paragraph{Amendment II}

The \bf{Second Amendment} safeguards the people's right to \it{"Bear Arms”}, meaning to own and carry weapons. Americans disagree about whether owning guns causes an increase or decrease in crime, as well as whether this amendment applies only to those who are members of the organized armed forces or National Guard. The \bf{District of Columbia v. Heller} case ruled that it protects individuals not connected to a state militia, at least for those residing in the nation’s capital.

U.S. law limits the \bf{Second Amendment}. It prohibits certain people from owning guns, like those found guilty of serious crimes. The \bf{Supreme Court} has said that it is acceptable for the federal government to make this law, but that it is up to the state governments to enforce it. Therefore, many states require background checks of people who wish to purchase guns to ensure they have not had this right taken away for a crime or other reason.

\begin{figure}[h]
    \begin{quote}
        \it{“A well regulated Militia, being necessary to the security of a free State, the right of the people to keep and bear Arms, shall not be infringed.”}
    \end{quote}
    \caption{Amendment II}
\end{figure}

\paragraph{Amendment III}

The \bf{Third Amendment} is a reaction to events that followed the \bf{French} and \bf{Indian War} that ended in $1763$. British soldiers remained in certain parts of the colonies after war’s end presumably to keep the peace. Some colonists had to provide for the troops’ basic needs, including letting them stay in their homes. They found this threatening and a violation of their rights. Therefore, when the framers wrote the Constitution, they wanted to make sure the new government could not do the same. However, this amendment does not prevent state or federal governments from using soldiers to restore order or assist civilians during an emergency.

\begin{figure}[h]
    \begin{quote}
        \it{“No Soldier shall, in time of peace be quartered in any house, without the consent of the Owner, nor in time of war, but in a manner to be prescribed by law.”}
    \end{quote}
    \caption{Amendment III}
\end{figure}

\paragraph{Amendment IV}

The \bf{Fourth Amendment} safeguards people from unwarranted:

\begin{itemize}
    \item Search
    \item Seizure
\end{itemize}

This means that people are protected from government officials searching their property and seizing, or taking, their property without a warrant. Police should not search a home or take anything from it without a judge's permission.

One example of a limit to this right is when a person is a suspect in a crime. A judge may give permission for officers to look for and take evidence when the owner is suspected of a crime. Another limit on this right would occur when officers see a crime in progress. They do not have to wait for a warrant to act in this situation. The officers that witness the crime have probable cause to act.

The \bf{Constitution} does not say directly that Americans have a right to privacy. However, the \bf{Supreme Court} has ruled that Americans generally have this right, supported by the Fourth Amendment.

What if police conduct a search without a warrant or probable cause and happen to find evidence of a crime, such as the sale of illegal drugs? The exclusionary rule states that evidence seized illegally cannot be used in a trial, though the courts have supported the good faith exception when the search error is minor. An example could be that police did have a warrant, but the judge wrote the wrong date on it.

\begin{figure}[h]
    \begin{quote}
        \it{“The right of the people to be secure in their persons, houses, papers, and effects, against unreasonable searches and seizures, shall not be violated, and no Warrants shall issue, but upon probable cause, supported by Oath or affirmation, and particularly describing the place to be searched, and the persons or things to be seized.”}
    \end{quote}
    \caption{Amendment IV}
\end{figure}

\paragraph{Amendment V}

The \bf{Fifth Amendment} protects several rights for those accused of crimes. It protects people from being tried for a crime without a grand jury’s review that there is sufficient evidence to suspect the person’s guilt. One exception here is a crime where the person was on active duty in the military at the time of the alleged offense. A person cannot go on trial for the same crime twice, referred to here as \bf{Double Jeopardy}. If authorities suspect the person of a similar crime later, they can go on trial again. Nor does a suspect have to testify if they do not wish, in order to protect themselves from the potential for self-incrimination. If a suspect goes to trial, it is up to the accusers to prove the person’s guilt. The suspect does not have to say anything (“pleading the Fifth”) as court procedure assumes innocence until proven otherwise. However, if a suspect does choose to speak, waving Fifth Amendment rights, the courts can use that speech as evidence.

No punishment of a person or taking of property can happen without the person having due process of law, meaning all the established laws related to investigating and judging the alleged crime have been followed. Should the government be able to take property away or jail someone while the investigation is ongoing? The answer has depended on the situation. Should the government be able to take property for a public reason, such as the need to build a highway? The amendment’s final line protects people from the government taking their property for the common good without paying them for it. Like the First Amendment, many controversies have arisen from this amendment.

\begin{figure}[h]
    \begin{quote}
        \it{“No person shall be held to answer for a capital, or otherwise infamous crime, unless on a presentment or indictment of a Grand Jury, except in cases arising in the land or naval forces, or in the Militia, when in actual service in time of War or public danger; nor shall any person be subject for the same offence to be twice put in jeopardy of life or limb; nor shall be compelled in any criminal case to be a witness against himself, nor be deprived of life, liberty, or property, without due process of law; nor shall private property be taken for public use, without just compensation.”}
    \end{quote}
    \caption{Amendment V}
\end{figure}

\paragraph{Amendment VI}

The \bf{Sixth Amendment} protects an accused person’s right to a jury trial that is \it{“Speedy and Public”} and in a court near the alleged crime. The purpose is to judge the person as fairly as possible with the fewest expenses and least time lost. Making a jury trial open to public view helps protect the rule of law, as due process is most likely to be followed when there is an audience. The jury should be \it{“Impartial”}, meaning jury members are not biased for or against the suspect. A suspect has the right to know details of the accusation and the evidence gathered, as well as to face the accusers in court. Suspects also have the right to have a lawyer represent them and present their own evidence and witnesses.

The \it{“Public”} part of this amendment has led to conflicts with the \bf{First Amendment} and privacy. Should news stations be able to broadcast evidence against the suspect before the court has selected all jury members? Should courts allow reporters in the courtroom when the victim of a crime is testifying? If jury selection goes on for months without producing enough unbiased members, should the trial move to another court? Judges have made decisions on related issues like these on a case-by-case basis.

\begin{figure}[h]
    \begin{quote}
        \it{“In all criminal prosecutions, the accused shall enjoy the right to a speedy and public trial, by an impartial jury of the State and district wherein the crime shall have been committed, which district shall have been previously ascertained by law, and to be informed of the nature and cause of the accusation; to be confronted with the witnesses against him; to have compulsory process for obtaining witnesses in his favor, and to have the Assistance of Counsel for his defence.”}
    \end{quote}
    \caption{Amendment VI}
\end{figure}

\paragraph{Amendment VII}

The \bf{Seventh Amendment} relates to civil court cases. A civil suit is a dispute between private citizens or groups of citizens, such as one person suing another for damaged property. The amendment gives people who seek settlement of the dispute through the courts the right to have a jury decide the case if they choose. It also prevents the courts from trying the same case again, like the protection for crime suspects in the Fifth Amendment. You may have seen cases like these on reality shows decided by a judge. The citizens on the show have waived their Seventh Amendment right to a jury trial in the case.

\begin{figure}[h]
    \begin{quote}
        \it{“In Suits at common law, where the value in controversy shall exceed twenty dollars, the right of trial by jury shall be preserved, and no fact tried by a jury, shall be otherwise re-examined in any Court of the United States, than according to the rules of the common law.”}
    \end{quote}
    \caption{Amendment VII}
\end{figure}

\paragraph{Amendment VIII}

The \bf{Eighth Amendment} protects people from \it{“Cruel or Unusual Punishment”}. Its purpose is to protect people from excessive punishment and bail requirements. Bail is the amount of money required to allow a suspect to be released from jail while awaiting trial. To guarantee that the suspect will show up in court, they can then have the funds returned. Yet who decides what is excessive? Is it cruel to sentence a child to prison who commits a serious crime? Can a judge set bail at one million dollars for a person convicted of writing a bad check?

The \bf{Eighth Amendment} is another source of intense debate among Americans. Judges make decisions and attempt to balance individual rights with the public good. For example, most people would agree that a person accused of murder should not have a right to bail at all, for fear the suspect may harm others while out of custody. People also disagree whether capital punishment, or the death penalty, is acceptable punishment for certain crimes. In general, Americans view that the more serious the crime, the higher expected bail and potential punishment.

\begin{figure}[h]
    \begin{quote}
        \it{“Excessive bail shall not be required, nor excessive fines imposed, nor cruel and unusual punishments inflicted.”}
    \end{quote}
    \caption{Amendment VIII}
\end{figure}

\paragraph{Amendment IX}

The \bf{Ninth Amendment} protects people from government laws that would unfairly restrict rights not named by the \bf{Bill of Rights}. Many of the Constitution’s framers worried that listing individual rights implied that people did not have any other rights. Therefore, as it stands by this amendment, if there is no law against a certain behavior, Americans can assume they have the right to it. The Constitution itself prevents ex post facto laws, so if a behavior is outlawed by the government after it was determined to be harmful to society, it cannot arrest people for actions occurring before the law passed.

The courts have often studied the amendment to judge cases that relate to an assumed right to privacy in lifestyle choices. The Constitution does not state directly that the people have a right to privacy. However, interpretations of the Fourth and Ninth Amendments have set precedent for this right where it does not conflict with the public good.

\begin{figure}[h]
    \begin{quote}
        \it{“The enumeration in the Constitution, of certain rights, shall not be construed to deny or disparage others retained by the people.”}
    \end{quote}
    \caption{Amendment IX}
\end{figure}

\paragraph{Amendment X}

The \bf{Tenth Amendment} protects the powers of the state governments, or \it{“States’ rights”}. It prevents the federal government from taking over responsibilities reserved to the states because the first three articles of the Constitution do not delegate them to the national government. The amendment helps preserve the limited central government the framers intended.

The federal government has expanded significantly since the writing of the Constitution. Americans debate today whether many of its programs and departments are reasonable due to the \it{“General Welfare”} and \it{“Necessary and Proper”} statements of the main \bf{Constitution}, or if the states have the right to challenge certain policies and programs under the \bf{Tenth Amendment}. For example, does a state have a right to refuse compliance with a national health care program, traditionally an area covered by state governments? The supremacy clause states that federal laws are supreme, but states have challenged federal programs they disagreed with through the courts.

\begin{figure}[h]
    \begin{quote}
        \it{“The powers not delegated to the United States by the Constitution, nor prohibited by it to the States, are reserved to the States respectively, or to the people.”}
    \end{quote}
    \caption{Amendment X}
\end{figure}

\subsubsection*{Perspectives on Balancing Liberty and Order}

Americans generally accept limitations to rights that are common sense for public safety. A school is an example of a place where laws may restrict individual rights to protect the common good. For example, students have the right to assemble under the \bf{First Amendment}. However, school officials may set rules around when, why, and how assemblies may occur. Restrictions help protect safety for everyone at the school. The \bf{Supreme Court} has ruled on various laws affecting rights not only in schools, but also in other public areas as well, like courthouses, parks, and roads.

\paragraph{Causes} Certain situations often lead people to support greater restrictions on individual rights on a temporary basis. Wars, natural disasters, and other emergency threats are the most common reasons for greater limits. For example, if soldiers are sent to war in another nation, the government may limit press coverage of their activities to protect them from enemy ambush. A bomb threat at a courthouse may lead officials to stop people from entering, exiting, or moving about the building until the threat is over. What happens when there is no clear end to a threat, such as a war that continues for years? A modern example is the efforts of the United States to prevent terrorism.

The \bf{USA PATRIOT} Act was legislation that expanded the powers of the federal government to prevent terrorist attacks from within, as well as outside, the country. \it{President Bush} signed it into law and launched a \it{"War on Terror"} in $2001$. Congress changed many parts of the act after its initial passage. In $2011$, \it{President Obama} signed an extension of certain parts of the \bf{USA PATRIOT Act} that affect civil liberties, including roving wiretaps. The \bf{USA PATRIOT Act} expired in $2015$. The \bf{USA FREEDOM Act} was similar legislation passed to replace it. In $2020$, Congress passed a reauthorization of the \bf{USA FREEDOM Act} which changed and extended some surveillance privileges.

\paragraph{USA PATRIOT Act Debate} Debate over the \bf{USA PATRIOT Act} has been steady since its passage. Questions of where the act stands on the goal of balancing liberty and security have ranged from the general—such as \it{"Does it violate civil liberties?"}—to the specific—such as \it{"Do roving wiretaps violate the Fourth Amendment?"}

\subsubsection*{Emergencies Justifying Restriction on Civil Liberties}

Americans often settle issues related to civil liberties in courtrooms, with lawyers and judges debating the intent of laws, the \bf{Constitution}, and the \bf{Bill of Rights}. People who take a legal statement at face value prefer a \it{"Strict"} interpretation of the law. Those who see legal statements as flexible in application agree with a \it{"Loose"} interpretation of the laws. No one is entirely \it{"Strict"} or \it{"Loose"} in all interpretations, however. Even \bf{Supreme Court} justices have shifted from one view to the other on similar situations because of different circumstances or laws involved.

\newpage
